\documentclass[a4paper]{article}
\usepackage{amsfonts}
\usepackage{amsmath}
\usepackage{amssymb}
\usepackage{graphicx}     

\begin{document}
  
\noindent \large Auteur: Victoria Mazur \\
\noindent \large Studierichting: Informatica\\
\noindent \large Oefeningenreeks 2B nr. 6.5\\

\medskip

\normalsize

$\alpha = \operatorname{Bgtan} \dfrac{1}{7}$ \\

$\beta = \operatorname{Bgtan} \dfrac{1}{3}$ \\

$\alpha + 2\beta$ \\

$\tan \alpha = \dfrac{1}{7}$ \\

Dubbele hoekformule: $\tan 2\theta = \dfrac{2 \tan \theta}{1 - \tan^2 \theta}$ \\

$\tan 2\beta = \dfrac{2 \tan \beta}{1 - \tan^2 \beta} = \dfrac{2 \cdot \dfrac{1}{3}}{1 - \left(\dfrac{1}{3}\right)^2}$ \\

$= \dfrac{\dfrac{2}{3}}{1 - \dfrac{1}{9}} = \dfrac{\dfrac{2}{3}}{\dfrac{9-1}{9}} = \dfrac{\dfrac{2}{3}}{\dfrac{8}{9}}$ \\

$= \dfrac{2}{3} \cdot \dfrac{9}{8} = \dfrac{18}{24} = \dfrac{3}{4}$ \\

$\Rightarrow \tan(\alpha + 2\beta)$ \\

Somformule: $\tan(A + B) = \dfrac{\tan A + \tan B}{1 - \tan A \cdot \tan B}$ \\

$= \dfrac{\tan \alpha + \tan 2\beta}{1 - \tan \alpha \cdot \tan 2\beta}$ \\

$= \dfrac{\dfrac{1}{7} + \dfrac{3}{4}}{1 - \dfrac{1}{7} \cdot \dfrac{3}{4}}$ \\

Teller: $\dfrac{1}{7} + \dfrac{3}{4} = \dfrac{4}{28} + \dfrac{21}{28} = \dfrac{25}{28}$ \\

Noemer: $1 - \dfrac{1}{7} \cdot \dfrac{3}{4} = 1 - \dfrac{3}{28} = \dfrac{28-3}{28} = \dfrac{25}{28}$ \\

$= \dfrac{\dfrac{25}{28}}{\dfrac{25}{28}} = 1$ \\

\begin{thebibliography}{999}
%\bibitem{a} Referentie
\end{thebibliography}
\end{document}